\section{TỔNG QUAN VỀ CÔNG TY CỔ PHẦN FPT}

\subsection{Thông tin pháp lý cơ bản}
Công ty Cổ phần FPT là một trong những doanh nghiệp công nghệ thông tin và viễn thông lớn nhất tại Việt Nam. Bảng \ref{tab:thong_tin_phap_ly} tóm tắt các thông tin pháp lý nền tảng của công ty:

\begin{table}[htbp]
\centering
\caption{Thông tin pháp lý cơ bản của Công ty Cổ phần FPT}
\label{tab:thong_tin_phap_ly}
\begin{tabularx}{\textwidth}{p{4.2cm} X}
\toprule
\textbf{Nội dung} & \textbf{Thông tin chi tiết} \\
\midrule
Tên đầy đủ & Công ty Cổ phần FPT \\
Tên tiếng Anh & FPT Corporation \\
Tên viết tắt & FPT \\
Mã chứng khoán & \textbf{FPT} \\
Sàn niêm yết & HOSE -- Sở Giao dịch Chứng khoán TP. Hồ Chí Minh \\
Ngày giao dịch đầu tiên & 13/12/2006 \\
Loại hình doanh nghiệp & Công ty Cổ phần \\
Mã số thuế & 0101248141 \\
Quyết định cổ phần hóa & Quyết định số 178/QĐ-TTg ngày 28/02/2002 của Thủ tướng Chính phủ \\
Trụ sở chính & Tòa nhà FPT Tower, số 10 Phạm Văn Bạch, P. Dịch Vọng Hậu, Q. Cầu Giấy, Hà Nội \\
Website chính thức & \url{https://fpt.com} \\
Người đại diện theo pháp luật & Ông Trương Gia Bình -- Chủ tịch Hội đồng Quản trị \\
Tổng Giám đốc / CEO & Ông Nguyễn Văn Khoa \\
Ngành nghề chính & Công nghệ thông tin, phần mềm, viễn thông, chuyển đổi số, giáo dục, bán dẫn \\
Số cổ phiếu đang lưu hành & 1.885.727.648 cổ phiếu (sau đợt phát hành cổ phiếu thưởng 10:1 năm 2026; giữa năm 2026 là 1.714.326.422 cổ phiếu) \\
Vốn điều lệ & 18.857.276.480.000 đồng (cập nhật cuối năm 2026) \\
\bottomrule
\end{tabularx}
\end{table}

\subsection{Giới thiệu chung về công ty}
FPT được thành lập ngày 13/9/1988 bởi 13 nhà khoa học trẻ giàu nhiệt huyết, dẫn dắt bởi ông Trương Gia Bình, với khát vọng xây dựng một tổ chức hùng mạnh dựa trên nền tảng khoa học công nghệ, đổi mới sáng tạo và phụng sự quốc gia. Khởi đầu với tên gọi \textbf{Công ty Công nghệ Thực phẩm} (Food Processing Technology), doanh nghiệp nhanh chóng nắm bắt xu thế công nghệ thế giới để chuyển trục sang công nghệ thông tin vào năm 1990.

Trải qua gần 4 thập kỷ, FPT đã phát triển từ một công ty khởi nghiệp quy mô nhỏ thành một tập đoàn công nghệ số 1 Việt Nam và có vị thế đáng kể trên bản đồ CNTT toàn cầu. FPT định vị là doanh nghiệp tiên phong về chuyển đổi số và dẫn đầu trong tư vấn, cung cấp, triển khai các dịch vụ, giải pháp công nghệ -- viễn thông toàn diện.

Hiện nay, FPT hoạt động đa dạng trong các lĩnh vực:
\begin{itemize}[noitemsep]
    \item Phát triển phần mềm và dịch vụ công nghệ thông tin phục vụ thị trường toàn cầu.
    \item Tích hợp hệ thống, tư vấn chiến lược và triển khai chuyển đổi số trong nước.
    \item Hạ tầng viễn thông, Internet băng rộng và truyền hình số.
    \item Điện toán đám mây (Cloud), Trí tuệ nhân tạo (AI), dữ liệu lớn, IoT, blockchain và công nghệ bán dẫn (Semiconductor).
    \item Giáo dục và đào tạo nguồn nhân lực công nghệ cao từ cấp phổ thông đến sau đại học.
    \item Bán lẻ và phân phối các sản phẩm, thiết bị công nghệ.
    \item Đầu tư mạo hiểm và ươm tạo các startup công nghệ.
\end{itemize}

\subsection{Lĩnh vực hoạt động, ngành nghề kinh doanh và sản phẩm tiêu biểu}

\subsubsection{Các mảng kinh doanh chính}
Hệ sinh thái kinh doanh của FPT được cơ cấu thành 10 mảng hoạt động chuyên sâu (Bảng \ref{tab:mang_kinh_doanh}):

\begin{table}[htbp]
\centering
\small
\caption{Các mảng kinh doanh chính trong hệ sinh thái FPT}
\label{tab:mang_kinh_doanh}
\begin{tabularx}{\textwidth}{p{3.8cm} X p{3.2cm}}
\toprule
\textbf{Mảng kinh doanh} & \textbf{Nội dung chính} & \textbf{Đơn vị tiêu biểu} \\
\midrule
Dịch vụ CNTT quốc tế & Phát triển phần mềm, outsourcing, AI, Cloud, kiểm thử, BPO & FPT Software \\
Giải pháp CNTT trong nước & Chuyển đổi số, chính phủ số, giải pháp tài chính, y tế & FPT IS \\
Điện toán đám mây & Cloud, hạ tầng số, trung tâm dữ liệu, nền tảng AI & FPT Smart Cloud \\
Chuyển đổi số & Tư vấn chiến lược số, số hóa quy trình doanh nghiệp & FPT Digital \\
Viễn thông & Dịch vụ Internet băng rộng, truyền hình số, cáp viễn thông & FPT Telecom \\
Giáo dục & Đào tạo kỹ sư phần mềm, AI, vi mạch và nhân lực số & FPT Education \\
Truyền thông số & Báo chí điện tử trực tuyến, giải pháp quảng cáo số & FPT Online (VnExpress) \\
Bán lẻ công nghệ & Chuỗi bán lẻ điện thoại, máy tính, thiết bị thông minh & FPT Retail (FPT Shop) \\
Phân phối công nghệ & Phân phối sản phẩm CNTT, điện tử tiêu dùng chính hãng & Synnex FPT \\
Đầu tư & Đầu tư chiến lược, M\&A, ươm tạo startup công nghệ & FPT Investment \\
\bottomrule
\end{tabularx}
\end{table}

\subsubsection{Các sản phẩm và nền tảng công nghệ tiêu biểu}
Hệ sinh thái sản phẩm \textit{``Made by FPT''} đã khẳng định uy tín trên thị trường trong và ngoài nước:
\begin{itemize}
    \item \textbf{FPT.AI:} Nền tảng Trí tuệ nhân tạo toàn diện cho doanh nghiệp (phát triển từ năm 2017), bao gồm chatbot, trợ lý ảo thông minh, xử lý ngôn ngữ tự nhiên (NLP), thị giác máy tính (Computer Vision) và tự động hóa hội thoại.
    \item \textbf{akaBot:} Nền tảng tự động hóa quy trình nghiệp vụ kết hợp RPA (Robotic Process Automation), OCR, IDP và Agentic AI; đã triển khai thành công cho hơn 4.500 khách hàng doanh nghiệp tại 21 quốc gia.
    \item \textbf{FPT Smart Cloud:} Hệ sinh thái điện toán đám mây và trung tâm dữ liệu thế hệ mới, đáp ứng tiêu chuẩn khắt khe về an toàn bảo mật.
    \item \textbf{VnExpress:} Báo điện tử tiếng Việt nhiều người đọc nhất, ra mắt từ năm 2001, giữ vị trí đầu ngành truyền thông số.
    \item \textbf{FPT Play:} Nền tảng truyền hình và giải trí số tương tác hàng đầu tại Việt Nam.
    \item \textbf{FPT Long Châu \& FPT Shop:} Chuỗi nhà thuốc bán lẻ hiện đại dẫn đầu cả nước và hệ thống bán lẻ công nghệ trực thuộc FPT Retail.
    \item \textbf{FPT Semiconductor:} Đơn vị thiết kế vi mạch, chip bán dẫn phục vụ thị trường vi mạch nhúng và thiết bị y tế, IoT.
    \item \textbf{Các giải pháp chính phủ số của FPT IS:} Hệ thống quản lý thuế, hải quan điện tử, hệ sinh thái tài chính công, bệnh viện thông minh và cổng dịch vụ công trực tuyến.
\end{itemize}

\subsection{Lịch sử hình thành và phát triển}
Hành trình gần 40 năm của FPT được chia thành các giai đoạn mang tính bước ngoặt:
\begin{enumerate}
    \item \textbf{1988--1990: Giai đoạn khởi đầu:} Thành lập ngày 13/9/1988 với tên gọi Công ty Công nghệ Thực phẩm. Ngày 27/10/1990, đổi tên thành \textit{The Corporation for Financing and Promoting Technology} (FPT), chính thức rẽ hướng sang CNTT.
    \item \textbf{1994--1999: Mở rộng công nghệ và bước đầu toàn cầu hóa:} Năm 1994, trở thành đại lý phân phối sản phẩm IBM tại Việt Nam; năm 1997 là 1 trong 4 đơn vị đầu tiên được cấp phép cung cấp dịch vụ Internet; năm 1999 tuyên bố chiến lược xuất khẩu phần mềm ra thế giới.
    \item \textbf{2000--2006: Cổ phần hóa và niêm yết chứng khoán:} Năm 2001 ra mắt báo VnExpress; năm 2002 hoàn tất cổ phần hóa; năm 2006 chính thức niêm yết cổ phiếu trên sàn HOSE và thành lập Trường Đại học FPT.
    \item \textbf{2008--2014: Tăng tốc hội nhập quốc tế:} Đổi tên thành Công ty Cổ phần FPT (2008); mở rộng bán lẻ FPT Shop (2012); thực hiện thương vụ M\&A quốc tế đầu tiên của ngành CNTT Việt Nam (mua lại RWE IT Slovakia năm 2014).
    \item \textbf{2017--2022: Mở rộng hệ sinh thái số:} Tăng vốn điều lệ liên tục, thúc đẩy mạnh mẽ các mảng Cloud, Big Data, AI và chuyển đổi số toàn diện.
    \item \textbf{2023--2026: Đột phá với AI, Bán dẫn và DC5:}
    \begin{itemize}
        \item Năm 2023--2024: Đầu tư mạnh vào công nghệ bán dẫn, phần mềm ô tô (Automotive) và AI.
        \item Năm 2024--2025: Tháng 4/2024, FPT công bố hợp tác cùng NVIDIA đầu tư khoảng 200 triệu USD xây dựng Tổ hợp Nhà máy Trí tuệ Nhân tạo (AI Factory) tại Việt Nam và ra mắt tại Nhật Bản tháng 11/2024. Tháng 9/2025, FPT ký hợp đồng chuyển đổi số tích hợp AI trị giá 256 triệu USD với một tập đoàn năng lượng châu Á.
        \item Năm 2025: Triển khai chiến lược \textbf{DC5 (Digital Conglomerate 5.0 / DC5-135)}, xây dựng tổ hợp doanh nghiệp số dựa trên 5 trụ cột (AI, Data, Định danh, Giao tiếp, Điểm chạm) cùng tầm nhìn phát triển 1 triệu nhân lực chuyển đổi số vào năm 2035.
        \item Năm 2026: Sau đợt phát hành ESOP giữa năm 2026 (đạt 1.714.326.422 cổ phiếu) và đợt phát hành cổ phiếu thưởng tỷ lệ 10:1 cuối năm 2026, vốn điều lệ đạt 18.857.276.480.000 đồng với 1.885.727.648 cổ phiếu lưu hành.
    \end{itemize}
\end{enumerate}

\subsection{Cơ cấu sở hữu và mạng lưới thành viên}

\subsubsection{Cơ cấu cổ đông lớn}
Theo số liệu trước đợt chia cổ phiếu thưởng 10:1 (tại ngày 30/06/2026, tổng số cổ phiếu lưu hành sau ESOP là 1.714.326.422 cổ phiếu; cuối năm 2025 là khoảng 1,70 tỷ cổ phiếu), ông Trương Gia Bình nắm giữ 117.347.966 cổ phiếu và SCIC nắm giữ 96.585.637 cổ phiếu. Sau khi hoàn tất phát hành cổ phiếu thưởng tỷ lệ 10:1 vào cuối năm 2026 nâng tổng lượng cổ phiếu lưu hành lên \textbf{1.885.727.648 cổ phiếu}, cơ cấu cổ đông lớn của FPT được tổng hợp qua Bảng \ref{tab:co_dong_lon}:

\begin{table}[htbp]
\centering
\caption{Cơ cấu cổ đông lớn tại Công ty Cổ phần FPT (cập nhật năm 2026)}
\label{tab:co_dong_lon}
\begin{tabularx}{\textwidth}{Xrr}
\toprule
\textbf{Tên cổ đông} & \textbf{Số cổ phiếu (sau chia 10:1)} & \textbf{Tỷ lệ sở hữu} \\
\midrule
Ông Trương Gia Bình & khoảng 129.082.762 cổ phiếu & khoảng \textbf{6,85\%} \\
Tổng công ty Đầu tư và Kinh doanh vốn Nhà nước (SCIC) & khoảng 106.244.200 cổ phiếu & khoảng \textbf{5,63\%} \\
Các cổ đông tổ chức, cá nhân trong và ngoài nước khác & khoảng 1.650.400.686 cổ phiếu & khoảng \textbf{87,52\%} \\
\midrule
\textbf{Tổng cộng} & \textbf{1.885.727.648 cổ phiếu} & \textbf{100,00\%} \\
\bottomrule
\end{tabularx}
\end{table}

Cơ cấu sở hữu của FPT có tính đại chúng rất cao, có sự tham gia của các quỹ đầu tư tài chính quốc tế uy tín, các cổ đông nội bộ và nhà đầu tư cá nhân, bảo đảm tính thanh khoản và hạn chế tình trạng thâu tóm chi phối độc đoán.

\subsubsection{Cơ cấu sở hữu công ty con và liên kết cấp 1}
Tập đoàn FPT nắm quyền sở hữu chi phối hoặc liên kết tại 10 đơn vị cấp 1 (Bảng \ref{tab:so_huu_cap_1}):

\begin{table}[htbp]
\centering
\caption{Cơ cấu sở hữu cấp 1 của FPT Corporation}
\label{tab:so_huu_cap_1}
\begin{tabularx}{\textwidth}{p{3.2cm} X r p{3.6cm}}
\toprule
\textbf{Công ty} & \textbf{Tên pháp nhân đầy đủ} & \textbf{Tỷ lệ sở hữu} & \textbf{Phân loại} \\
\midrule
FPT Software & Công ty TNHH Phần mềm FPT & 100\% & Công ty con \\
FPT IS & Công ty TNHH FPT IS & 100\% & Công ty con \\
FPT Smart Cloud & Công ty TNHH FPT Smart Cloud & 100\% & Công ty con \\
FPT Digital & Công ty TNHH FPT Digital & 100\% & Công ty con \\
FPT Education & Công ty TNHH Giáo dục FPT & 100\% & Công ty con \\
FPT Investment & Công ty TNHH Đầu tư FPT & 100\% & Công ty con \\
FPT Online & Công ty Cổ phần Dịch vụ Trực tuyến FPT & 49,52\%* & Công ty con \\
FPT Telecom & Công ty Cổ phần Viễn thông FPT & 45,66\%** & Công ty liên kết (từ 2026) \\
FPT Retail & Công ty Cổ phần Bán lẻ Kỹ thuật số FPT & 46,54\% & Công ty liên kết \\
Synnex FPT & Công ty Cổ phần Synnex FPT & 48,00\% & Công ty liên kết \\
\bottomrule
\end{tabularx}
\end{table}

\noindent\footnotesize{\textit{Ghi chú:} (*) FPT nắm trực tiếp 49,52\% vốn tại FPT Online nhưng thông qua các đơn vị trong hệ sinh thái đạt tổng quyền biểu quyết trên 80\%, nên được hợp nhất là công ty con. (**) Tháng 7/2025, SCIC đã chuyển giao 50,17\% vốn nhà nước tại FPT Telecom sang Bộ Công an (và tháng 10/2026 bàn giao về Tổng công ty GTINX); từ đầu năm 2026, FPT chuyển FPT Telecom từ công ty con sang hạch toán công ty liên kết.}\normalsize

\subsubsection{Mạng lưới công ty con và công ty liên kết cấp 2}
Nhằm mở rộng thị trường và tối ưu hóa chuyên môn hóa, các công ty cấp 1 thiết lập các pháp nhân cấp 2:
\begin{itemize}
    \item \textbf{Dưới FPT Software:} Mạng lưới các công ty con toàn cầu đảm nhận các thị trường chiến lược bao gồm \textit{FPT Software Japan} (Nhật Bản), \textit{FPT Software Americas} (Hoa Kỳ), \textit{FPT Software Europe} (châu Âu), \textit{FPT Software Asia Pacific} và \textit{FPT Software India}; công ty công nghệ bán dẫn \textit{FPT Semiconductor}; nắm giữ 10\% cổ phần tại công ty dịch vụ BPO quốc tế \textit{Daythree Digital Berhad}.
    \item \textbf{Dưới FPT IS:} Các đơn vị thành viên theo chuyên môn như \textit{FPT-IS-BANK} (chuyên giải pháp phần mềm lõi ngân hàng) và các đơn vị chính phủ số, y tế số.
    \item \textbf{Dưới FPT Retail:} Quản lý chuỗi nhà thuốc \textit{FPT Long Châu} (đơn vị động lực tăng trưởng doanh thu vượt bậc) và chuỗi cửa hàng bán lẻ công nghệ \textit{FPT Shop / F.Studio by FPT}.
    \item \textbf{Dưới FPT Education:} Hệ sinh thái gồm \textit{Trường Đại học FPT (FPT University)}, hệ thống trường phổ thông liên cấp \textit{FPT School}, các viện đào tạo quốc tế và phân hiệu chuyên môn.
\end{itemize}
